% vim:foldmethod=marker:foldmarker=(((,))):foldlevel=0

\usepackage{thmtools}
\usepackage[framemethod=TikZ]{mdframed}
\mdfsetup{skipabove=.5em, skipbelow=0em}

%: Theorem styles (((
\theoremstyle{definition}

\declaretheoremstyle[
    headfont=\bfseries\sffamily\color{ForestGreen!70!black}, bodyfont=\normalfont,
    mdframed={
        linewidth=2pt,
        rightline=false, topline=false, bottomline=false,
        linecolor=ForestGreen, backgroundcolor=ForestGreen!5,
    }
]{thmgreenbox}

\declaretheoremstyle[
    headfont=\bfseries\sffamily\color{NavyBlue!70!black}, bodyfont=\normalfont,
    mdframed={
        linewidth=2pt,
        rightline=false, topline=false, bottomline=false,
        linecolor=NavyBlue, backgroundcolor=NavyBlue!5,
    }
]{thmbluebox}

\declaretheoremstyle[
    headfont=\bfseries\sffamily\color{NavyBlue!70!black}, bodyfont=\normalfont,
    mdframed={
        linewidth=2pt,
        rightline=false, topline=false, bottomline=false,
        linecolor=NavyBlue
    }
]{thmblueline}

\declaretheoremstyle[
    headfont=\bfseries\sffamily\color{RawSienna!70!black}, bodyfont=\normalfont,
    mdframed={
        linewidth=2pt,
        rightline=false, topline=false, bottomline=false,
        linecolor=RawSienna, backgroundcolor=RawSienna!5,
    }
]{thmredbox}

\declaretheoremstyle[
    headfont=\bfseries\sffamily\color{RawSienna!70!black}, bodyfont=\normalfont,
    mdframed={
        linewidth=2pt,
        rightline=false, topline=false, bottomline=false,
        linecolor=RawSienna, backgroundcolor=RawSienna!2.5
    }
]{thmredline}

\declaretheoremstyle[
    headfont=\bfseries\sffamily\color{RawSienna!70!black}, bodyfont=\normalfont,
    numbered=no,
    mdframed={
        linewidth=2pt,
        rightline=false, topline=false, bottomline=false,
        linecolor=RawSienna, backgroundcolor=RawSienna!1,
    },
    qed=\qedsymbol
]{thmproofbox}

\newcommand\chevron{\square\hspace{-0.82em}\raisebox{0.41em}{\resizebox{0.85em}{0.25em}{\triangledown}}}

\declaretheoremstyle[
    headfont=\bfseries\sffamily,
    bodyfont=\normalfont,
    numbered=yes,
    mdframed={
        linewidth=2pt,
        rightline=false, topline=false, bottomline=false,
        linecolor=Black, backgroundcolor=Black!1,
    },
    qed=\chevron
]{thproblembox}

%: )))

%: Theorems (((

%: Simple theorems (((

\declaretheorem[style=thmgreenbox, name=Определение]{definition}

\declaretheorem[style=thmbluebox, numbered=no, name=Пример]{eg}
\declaretheorem[style=thmbluebox, numbered=no, name=Примеры]{egs}
\declaretheorem[style=thmbluebox, numbered=no, name=Свойство]{property}
\declaretheorem[style=thmbluebox, numbered=no, name=Свойства]{properties}
\declaretheorem[style=thmbluebox, numbered=no, name=Алгоритм]{algorithm}

\declaretheorem[style=thmredbox, name=Теорема]{theorem}
\declaretheorem[style=thmredbox, name=Лемма]{lemma}
\declaretheorem[style=thmredbox, name=Утверждение]{statement}
\declaretheorem[style=thmredbox, name=Аксиома]{axiom}
\declaretheorem[style=thmredbox, numbered=no, name=Утверждения]{statements}
\declaretheorem[style=thmredbox, numbered=no, name=Следствие]{implication}
\declaretheorem[style=thmredbox, numbered=no, name=Вывод]{corollary}

\declaretheorem[style=thmblueline, numbered=no, name=Замечание]{remark}
\declaretheorem[style=thmblueline, numbered=no, name=Примечание]{note}

\declaretheorem[style=thproblembox, name=Задача]{problem}

\newtheorem*{notation}{Обозначение}
\newtheorem*{remind}{Напоминание}

%: )))

%: Enumerated theorems (((

\newenvironment{props}[1][ ]{\begin{properties}[{#1}]\hfill\begin{enumerate}}{\end{enumerate}\end{properties}}
\newenvironment{stmts}[1][ ]{\begin{statements}[{#1}]\hfill\begin{enumerate}}{\end{enumerate}\end{statements}}
\newenvironment{exmpls}[1][ ]{\begin{egs}[{#1}]\hfill\begin{enumerate}}{\end{enumerate}\end{egs}}
\newenvironment{algo}[1][ ]{\begin{algorithm}[{#1}]\hfill\begin{enumerate}}{\end{enumerate}\end{algorithm}}

\newenvironment{iproof}[1][ ]{\begin{proof}[{#1}]\hfill\begin{itemize}}{\end{itemize}\end{proof}}
\newenvironment{eproof}[1][ ]{\begin{proof}[{#1}]\hfill\begin{enumerate}}{\end{enumerate}\end{proof}}

%: )))

%: Raised theorems (((

\declaretheorem[style=thmredline, numbered=no, name=Другая формулировка]{restatebase}
\newenvironment{restate}{\vspace{-.5em}\begin{restatebase}}{\end{restatebase}}

\declaretheorem[style=thmproofbox, name=Доказательство]{replacementproof}
% proof, moved 0.5em up
\renewenvironment{proof}{\vspace{-.5em}\begin{replacementproof}}{\end{replacementproof}}

\declaretheorem[style=thmproofbox, name=Без доказательства]{noproofbase}
% Reasons go inside
\newenvironment{noproof}{\vspace{-.5em}\begin{noproofbase}}{\end{noproofbase}}

%: )))

%: Theorem-like environments (((

% "undefined theorem" -- theorem-like environment with custom header
\newmdenv[linewidth=2pt, rightline=false, topline=false, bottomline=false, linecolor=black, backgroundcolor=black!2]{undefinedtheorem}
\newenvironment{undefthm}[1]{\begin{undefinedtheorem}\textbf{#1 }}{\end{undefinedtheorem}}

% markdown quote
\newmdenv[linewidth=2pt, rightline=false, topline=false, bottomline=false, linecolor=black!25,
backgroundcolor=black!2]{mquote}

%: )))

%: )))
